\documentclass[11pt]{article}

\usepackage[margin=1in]{geometry}
\usepackage{amsmath,amssymb}
\usepackage{booktabs}
\usepackage{graphicx}
\usepackage{microtype}
\usepackage[numbers,sort&compress]{natbib}
\usepackage[hidelinks]{hyperref}
\usepackage{url}
\usepackage{listings}
\usepackage{xcolor}

\lstset{
  basicstyle=\ttfamily\footnotesize,
  breaklines=true,
  frame=single,
  framesep=4pt,
  rulecolor=\color{gray!40},
  columns=fullflexible,
}

\title{\textbf{Payload Bytes, Not Proxies:}\\
Diagnosing a Silent Data-Pipeline Failure in a\\
Transformer-Based Network Intrusion Detector}

\author{%
  Laalini Bhogadi\\
  Anvyon LLP\\
  \texttt{laalinib17@gmail.com}
}

\date{}

\begin{document}
\maketitle

\begin{abstract}
We report a diagnostic study of a 115.6M-parameter two-stage Transformer for
network intrusion detection that reached 75\% accuracy and emitted an anomaly
score confined to the band $[0.80, 0.85]$ regardless of input. The failure was
not architectural. Payload tensors---the input to the convolutional encoder that
holds 99\% of the model's parameters---were synthesized by hashing tabular flow
records. We show this destroys learnable structure: two flows differing only in
duration (0.100\,s vs.\ 0.101\,s) receive payload tensors with correlation
$-0.101$ and zero shared bytes out of 128. A gradient-boosted-tree probe
quantifies the resulting information ceiling at ROC-AUC $0.663$ on the features
the model actually received, against $0.9835$ on the full feature set. The
pipeline additionally exhibited a train/serve mismatch: training consumed hashed
proxies while inference consumed real packet bytes.

Replacing the proxy pipeline with real per-packet payload bytes from UNSW-NB15,
and unifying tokenization across training and inference, allows a
0.8M-parameter model to reach ROC-AUC $0.9880$---above the tabular
ceiling---while a 306.1M model reaches $0.9458$ after a single epoch. We report
three further findings of practical value: (i) the dispersion of predicted
scores detects representation collapse where the loss curve does not; (ii) on
this dataset a 0.8M model outperforms 306M and 517M variants, indicating the
task is data- rather than capacity-limited; and (iii) some XLA GPU versions
expand \texttt{float32} compilation of this architecture by a factor of
${\sim}180$, requesting 193\,GiB for a graph whose compiled footprint is
2.8\,GB, which we mitigate by mixed-precision compute and a pre-run compilation
probe. All code, logs, and the trained checkpoint are released.
\end{abstract}

\section{Introduction}

Machine learning for network intrusion detection has a well-documented
reproducibility problem. \citet{sommer2010outside} argued that the domain's
particular structure---extreme class imbalance, high cost of false positives,
and enormous variability in ``normal'' traffic---makes it unusually easy to
report strong results that do not transfer. \citet{arp2022dodos} catalogued the
specific pitfalls that recur in security ML papers, among them sampling bias,
inappropriate baselines, and \emph{lab-only evaluation} in which a model is
trained and tested under conditions it will never encounter in deployment.

This paper is an empirical case study of a failure mode adjacent to those
catalogues but, to our knowledge, not explicitly named: a data pipeline that
silently supplies an input carrying no learnable signal, while every conventional
training diagnostic---loss decreasing, no NaNs, checkpoints saving, evaluation
running---continues to look healthy.

The system under study is a two-stage Transformer over network flows, trained on
UNSW-NB15 \citep{moustafa2015unsw}. Its released checkpoint achieved 75\%
accuracy on a binary anomaly task and, when deployed against packet captures,
returned anomaly probabilities confined to $[0.80, 0.85]$ for every flow---
attack and benign alike. The natural reading is that the model is
under-trained or the architecture is unsuited to the task. Neither was correct.

\paragraph{Contributions.}
\begin{enumerate}
  \item We diagnose the failure to a data-pipeline defect and quantify it: the
  payload tensor, which feeds the encoder holding 99\% of parameters, was a
  deterministic hash of tabular fields and therefore carries no metric structure
  (Section~\ref{sec:diagnosis}).
  \item We introduce a simple, model-agnostic \emph{ceiling probe}: fit a
  gradient-boosted-tree classifier separately on the full feature set and on the
  features the model actually receives. The gap bounds what any architecture can
  achieve on the given input (Section~\ref{sec:ceiling}).
  \item We identify a train/serve tokenization mismatch in which the deployed
  model was served a distribution it had never trained on, and eliminate it by
  construction with a shared tokenizer and an equivalence test
  (Section~\ref{sec:skew}).
  \item We report that \emph{score dispersion} is a more reliable collapse
  detector than loss, and recommend logging it as standard practice
  (Section~\ref{sec:dispersion}).
  \item We document a systems failure in which \texttt{float32} compilation of
  this architecture expands ${\sim}180\times$ on some XLA GPU versions, and
  release a preflight tool that catches such cases in about a minute
  (Section~\ref{sec:systems}).
\end{enumerate}

\section{System Under Study}
\label{sec:system}

\subsection{Architecture}

The model processes a network flow as a sequence of $T$ packets, each truncated
to $L$ payload bytes. Writing $B$ for batch size:

\begin{enumerate}
  \item \textbf{Payload encoder.} Byte embedding over a vocabulary of $259$
  (256 byte values plus \textsc{pad}, \textsc{bos}, \textsc{eos}), learned
  positional embeddings, two 1-D convolutions with kernel width 3 and GELU
  activation, then mean-pooling over the byte axis. Input
  $[B, T, L] \to [B, T, E]$.
  \item \textbf{Header encoder.} Protocol and port-bucket embeddings
  concatenated with a projection of scalar header fields.
  \item \textbf{Packet stage.} $N_p$ pre-norm Transformer blocks
  \citep{vaswani2017attention} over the $T$ packet representations.
  \item \textbf{Protocol adapter.} A protocol embedding plus a bottleneck
  adapter.
  \item \textbf{Memory module.} Slot attention \citep{locatello2020object} pools
  packets into $S$ slots, whose mean is broadcast back as a session state.
  \item \textbf{Session stage.} $N_s$ further Transformer blocks, mean-pooled,
  feeding anomaly, traffic-class, and byte-decode heads.
\end{enumerate}

We use $L = 128$, $T = 16$, and $S = 8$ throughout. Three configurations are
studied, at $115{,}643{,}015$, $306{,}103{,}367$, and $517{,}221{,}903$
parameters. In all of them the Transformer blocks account for 98--99\% of
parameters, and the payload encoder is the only component operating on
$T \times L = 2048$ positions per sample rather than $T = 16$.

\subsection{The original data pipeline}

UNSW-NB15's widely used partitioned CSVs contain flow \emph{statistics}---
duration, byte and packet counts, TTLs, load, jitter, connection counters---but
no packet bytes. To supply the byte-level input the payload encoder requires,
the original pipeline synthesized one:

\begin{lstlisting}[language=Python]
def hash_bytes_from_fields(values, length):
    s = "|".join(str(v) for v in values).encode("utf-8")
    seed = int.from_bytes(hashlib.blake2b(s, digest_size=8).digest(),
                          "little") % (2 ** 32)
    rng = np.random.default_rng(seed)
    return rng.integers(0, 256, size=length, dtype=np.uint8).tobytes()
\end{lstlisting}

Each flow's fields are hashed to a seed, and a pseudo-random byte string is
drawn from it. The mapping is deterministic and reproducible, which is presumably
why it passed review.

\section{Diagnosis}
\label{sec:diagnosis}

\subsection{Symptom}

Training logs show validation accuracy identical across consecutive epochs
($0.7471$, $0.7471$), and a released checkpoint scoring $0.75$ accuracy with
macro-$F_1$ $0.7468$. At inference the model returned scores in
$[0.8008, 0.8524]$ across all evaluated captures. A frozen score band with a
decreasing loss is the signature of a model that has learned the class prior and
nothing else.

\subsection{Hash proxies carry no learnable structure}

A hash is designed to destroy locality: inputs that are close in the source
space map to outputs that are uncorrelated. Neural networks require the
opposite. We measured the effect directly on two flows differing only in
duration ($0.100$\,s vs.\ $0.101$\,s), holding all other fields fixed:

\begin{table}[h]
\centering
\begin{tabular}{lr}
\toprule
Quantity & Value \\
\midrule
Pearson correlation of the two payload tensors & $-0.101$ \\
Identical bytes (of 128) & $0$ \\
Correlation of a row with itself & $1.000$ \\
\bottomrule
\end{tabular}
\caption{Near-identical flows receive uncorrelated payload tensors. The mapping
is reproducible (correlation 1.0 on repeat), so the input is memorizable but not
generalizable---a one-time pad keyed by the record.}
\label{tab:hash}
\end{table}

A second, independent defect compounded this. The loader reads columns
\texttt{sport} and \texttt{dsport}, which do not exist in the UNSW-NB15
partitioned CSVs. Pandas' \texttt{Series.get} returns the supplied default
silently, so both port buckets were constant for every record. The header vector
reaching the model was therefore
$[\texttt{proto\_id}, 0, 0, \min(n_{\text{pkts}}, 16), 0]$: of 42 available
fields, exactly two carried information.

\subsection{Quantifying the information ceiling}
\label{sec:ceiling}

To separate ``the model is bad'' from ``the input is impoverished,'' we fit a
gradient-boosted-tree classifier \citep{friedman2001greedy} twice: once on the
full tabular feature set, and once on only the two features the Transformer
actually received. Training used the 175{,}341-record partition, evaluation the
82{,}332-record partition.

\begin{table}[h]
\centering
\begin{tabular}{lccc}
\toprule
Input & Accuracy & Macro $F_1$ & ROC-AUC \\
\midrule
All 42 tabular features & $0.8728$ & $0.8672$ & $\mathbf{0.9835}$ \\
Only $[\texttt{proto}, \min(n_{\text{pkts}},16)]$ & $0.6630$ & $0.5864$ & --- \\
Majority class & $0.5506$ & --- & --- \\
\bottomrule
\end{tabular}
\caption{The ceiling probe. The features reaching the model support at most
$\approx 0.66$ accuracy; the discarded features support ROC-AUC $0.9835$. The
reported model accuracy of $0.75$ sits in this gap.}
\label{tab:ceiling}
\end{table}

The probe costs seconds to run and bounds what \emph{any} architecture can
achieve on a given input. We recommend it as a routine check before attributing
poor performance to model capacity or optimization. Scaling the model under
these conditions cannot help: additional capacity buys memorization of
noise, not generalization.

\subsection{Train/serve tokenization mismatch}
\label{sec:skew}

The inference path parsed real \texttt{.pcap} files, extracted true payload
bytes, and tokenized them as $\texttt{BYTE\_OFFSET} + b$. Training used hashed
proxies. The deployed model was therefore evaluated on a distribution it had
never observed during training---a mismatch sufficient on its own to render the
scores meaningless, independent of the training defect.

This class of error is difficult to catch by inspection because both paths are
individually reasonable; only their conjunction is wrong. Our corrected
implementation exposes one tokenizer used by both paths and asserts their
equivalence in a unit test over payload lengths
$\{0, 1, 7, 64, 128\}$, including the case of a genuine \texttt{0x00} payload
byte, which must map to \texttt{BYTE\_OFFSET} rather than to padding.

\section{Corrected Pipeline}
\label{sec:corrected}

We replace synthesized bytes with real ones, drawn from the \texttt{Payload-Bytes}
subset of a packet-level redistribution of UNSW-NB15
\citep{pahalavan2023nids}, which provides per-packet payload bytes extracted
from the original captures together with flow identifiers and labels.

Packets are grouped by \texttt{flow\_id}, ordered by \texttt{packet\_id}, and the
first $T = 16$ retained; the first $L = 128$ bytes of each are tokenized. Because
Parquet is columnar, reading $130$ of $1485$ columns costs approximately
$130/1485$ of each file, so the pipeline streams directly from remote storage
without materializing the 47.9\,GB source. Flows, never packets, are split
between training and validation, so packets from one flow cannot straddle the
boundary.

Two further changes follow from properties of the corrected data:

\paragraph{Class weighting.} The true attack rate at flow level is $3.6\%$.
Unweighted cross-entropy reaches a low loss by predicting ``normal''
universally---a collapsed detector with a healthy loss curve. We apply
inverse-frequency weights, computed per split.

\paragraph{Fixed protocol vocabulary.} The tabular loader derived its protocol
vocabulary from the observed data, while the inference path maps IANA protocol
numbers modulo a fixed vocabulary size. We fix the vocabulary at 132 in both, so
that a protocol embedding denotes the same protocol at training and serving time.

\section{Experiments}
\label{sec:experiments}

\subsection{Setup}

Preprocessing yields $1{,}593{,}481$ flows from 18 source files
($1{,}434{,}132$ train / $159{,}349$ validation), reproduced identically across
three independent runs in two datacenter regions. Training uses AdamW
\citep{loshchilov2019adamw} with a cosine schedule and linear warmup, gradient
clipping at global norm 1.0, and batch size 256 on a single NVIDIA H100 80GB.
We report ROC-AUC as the primary metric: at a $3.6\%$ positive rate, accuracy is
dominated by the prior, and the class weighting deliberately shifts the optimal
threshold away from $0.5$.

\subsection{The pipeline fix recovers the signal}

A deliberately small model (0.8M parameters, $d_{\text{model}} = 128$) trained on
a single shard (89{,}337 flows) establishes that the corrected pipeline carries
signal:

\begin{table}[h]
\centering
\begin{tabular}{lccccc}
\toprule
Epoch & Loss & Accuracy & Macro $F_1$ & ROC-AUC & Score $p_{5}$--$p_{95}$ \\
\midrule
1 & $0.3497$ & $0.9400$ & $0.8220$ & $0.9689$ & $0.026$--$0.970$ \\
2 & $0.1901$ & $0.9609$ & $0.8732$ & $0.9736$ & $0.012$--$0.967$ \\
3 & $0.1488$ & $0.9665$ & $0.8883$ & $0.9843$ & $0.004$--$0.950$ \\
4 & $0.1188$ & $0.9710$ & $0.9006$ & $0.9853$ & $0.005$--$0.979$ \\
5 & $0.0995$ & $0.9457$ & $0.8390$ & $0.9858$ & $0.003$--$0.974$ \\
6 & $0.0865$ & $0.9715$ & $0.9033$ & $\mathbf{0.9880}$ & $0.003$--$0.976$ \\
\bottomrule
\end{tabular}
\caption{A 0.8M-parameter model on $1/18$ of the data exceeds the $0.9835$
tabular ceiling of Table~\ref{tab:ceiling}. The architecture was never the
limiting factor.}
\label{tab:small}
\end{table}

Epoch 5 is instructive: accuracy falls ($0.9710 \to 0.9457$) while ROC-AUC
\emph{rises} ($0.9853 \to 0.9858$). The ranking improved; only the fixed $0.5$
threshold became less apt. Practitioners monitoring accuracy would have read
this as a regression.

\subsection{Scale}

\begin{table}[h]
\centering
\begin{tabular}{lrrrc}
\toprule
Model & Parameters & Data & Epochs & ROC-AUC \\
\midrule
Protogrok (tiny) & $0.8$M & 1 shard & 6 & $\mathbf{0.9880}$ \\
GBDT, full tabular features & --- & full & --- & $0.9835$ \\
Protogrok (large) & $306.1$M & 18 shards & 1 & $0.9458$ \\
Protogrok (XL), $\eta = 6\!\times\!10^{-4}$ & $517.2$M & 18 shards & 3 & $0.8385$ \\
Protogrok, hashed proxies & $115.6$M & full CSV & 2 & $\approx 0.66$ ceiling \\
\bottomrule
\end{tabular}
\caption{Capacity does not help on this dataset. The smallest model is the best.}
\label{tab:scale}
\end{table}

The 517M configuration at $\eta = 6\times10^{-4}$ oscillated
($0.8348 \to 0.8048 \to 0.8385$ over three epochs) with loss decreasing only
marginally---a learning rate too large for the model width. Halving it to
$2.5\times10^{-4}$ for the 306M configuration moved first-epoch ROC-AUC from
$0.8348$ to $0.9458$ and widened the score range from $[0.301, 0.926]$ to
$[0.007, 0.969]$.

We emphasize the negative result in Table~\ref{tab:scale}. With
$1.43\times10^{6}$ training flows and 42 underlying features, a
$3\times10^{8}$-parameter model has roughly 200 parameters per training example.
The task is data-limited, not capacity-limited, and the original instinct to
scale to 500M parameters was the wrong axis of improvement.

\subsection{Score dispersion detects collapse}
\label{sec:dispersion}

Across every run, the interval between the 5th and 95th percentiles of the
predicted anomaly probability separated working from collapsed models more
reliably than the loss:

\begin{table}[h]
\centering
\begin{tabular}{llc}
\toprule
Condition & Score $p_5$--$p_{95}$ & Status \\
\midrule
Hashed proxies, deployed & $0.801$--$0.852$ & collapsed \\
517M, $\eta$ too large & $0.301$--$0.926$ & partially collapsed \\
306M, corrected $\eta$ & $0.007$--$0.969$ & healthy \\
0.8M, corrected pipeline & $0.003$--$0.976$ & healthy \\
\bottomrule
\end{tabular}
\caption{Score dispersion as a collapse diagnostic. In every collapsed case the
loss curve appeared unremarkable.}
\label{tab:dispersion}
\end{table}

The measure costs two percentile computations per evaluation and requires no
labels. We recommend logging it alongside loss as standard practice for
imbalanced detection tasks.

\subsection{Out-of-distribution saturation}

The released 306M checkpoint ranks well on held-out UNSW-NB15
(ROC-AUC $0.9458$) but saturates on captures unlike its training distribution:

\begin{table}[h]
\centering
\begin{tabular}{lccccc}
\toprule
Capture & Min & Median & Max & Flagged \\
\midrule
Crafted attack traffic & $0.9627$ & $0.9714$ & $0.9933$ & $41/41$ \\
Crafted benign + scan & $0.9627$ & $0.9715$ & $0.9933$ & $34/34$ \\
Live capture & $0.6797$ & $0.9570$ & $0.9922$ & $24/24$ \\
\bottomrule
\end{tabular}
\caption{On out-of-distribution captures the attack and benign score
distributions are indistinguishable. This is precisely the generalization gap
\citet{sommer2010outside} warned of, and it is invisible to in-distribution
validation.}
\label{tab:ood}
\end{table}

We report this because it is the outcome an in-distribution AUC of $0.9458$ does
not predict, and because the checkpoint is publicly released: a user who
evaluated it only on held-out UNSW-NB15 would substantially overestimate its
utility.

\section{Systems Findings}
\label{sec:systems}

\subsection{A \texttt{float32} compilation blowup}

Three training attempts differing in one variable:

\begin{table}[h]
\centering
\begin{tabular}{llll}
\toprule
Configuration & Batch & Compute dtype & Outcome \\
\midrule
$115.6$M & 32 & \texttt{float32} & OOM, $193.78$\,GiB requested \\
$306.1$M & 256 & \texttt{float32} & OOM, $2.02$\,TiB requested \\
$517.2$M & 256 & \texttt{bfloat16} & compiled at $15.79$\,GB \\
\bottomrule
\end{tabular}
\caption{Identical code; only the compute dtype differs. The larger model in
\texttt{bfloat16} compiles comfortably where a model $4.5\times$ smaller in
\texttt{float32} requests more memory than any accelerator provides.}
\label{tab:blowup}
\end{table}

Compiling the same training step for the CPU backend reports a total footprint
of $2.8$\,GB at batch 32, scaling linearly with batch size, establishing that the
computation does not require the requested memory. The rematerialization pass
reported it could not reduce the requirement at all
(``$193.78$\,GiB down from $193.78$\,GiB''), indicating a single unshrinkable
buffer rather than accumulated activations. We did not isolate the responsible
operator; we report the correlation because it is reproducible, severe, and
silently attributed to model size by anyone who has not measured otherwise.

The practical mitigation is mixed precision---\texttt{bfloat16} compute with
\texttt{float32} parameters---which is independently desirable on tensor-core
hardware.

\subsection{Preflight compilation probes}

Because the failure above manifests only after preprocessing and only on the
target device, we release a tool that compiles the training step across a matrix
of configuration variants and reports each one's measured footprint before any
long run. On an H100 80GB at batch 256 for the 517M configuration:

\begin{table}[h]
\centering
\begin{tabular}{llrrr}
\toprule
Conv. implementation & Remat & Temp (GB) & Args (GB) & Total (GB) \\
\midrule
\texttt{conv} & no & $9.59$ & $6.21$ & $15.79$ \\
\texttt{conv} & yes & $4.72$ & $6.21$ & $\mathbf{10.93}$ \\
\texttt{matmul} & no & $12.27$ & $6.21$ & $18.48$ \\
\texttt{matmul} & yes & $8.72$ & $6.21$ & $14.93$ \\
\bottomrule
\end{tabular}
\caption{Preflight output. Gradient checkpointing \citep{chen2016training}
halves activation memory on GPU; on the CPU backend the same setting changes
nothing ($20.13 \to 20.41$\,GB), so the probe must run on the target device.}
\label{tab:preflight}
\end{table}

The probe takes approximately one minute and would have prevented several hours
of failed accelerator time in our own experiments.

\section{Limitations}

The released checkpoint is trained for a single epoch; the run was terminated by
a full disk rather than by convergence, and validation AUC was still improving.
Its threshold is uncalibrated, and Table~\ref{tab:ood} documents that it does not
separate attack from benign traffic on out-of-distribution captures. We retain
only the first 16 packets per flow and 128 bytes per packet, whereas flows in
this corpus average 28 packets and median payloads of approximately 1300 bytes;
most of each payload is therefore unobserved. Rare attack classes are extremely
sparse (155 worms, 401 backdoor, 435 analysis out of $1.59\times10^6$ flows), so
the multi-class head is not evaluated here. Finally, all results are on a single
dataset; UNSW-NB15 is synthetic in its attack generation, and the
generalization concerns raised by \citet{sommer2010outside} apply in full.

\section{Conclusion}

The model examined here appeared to be a capacity or optimization problem and was
in fact a data problem. The diagnostic sequence that resolved it---measure
whether the input carries structure, bound the achievable performance with a
cheap non-neural probe, verify that training and serving consume the same
representation, and monitor score dispersion rather than loss alone---is
inexpensive and generalizes beyond this system. We note that our own initial
instinct, to scale from 115M to 500M parameters, would have made the situation
worse while appearing to be progress; the smallest model we trained remains the
best-performing.

\section*{Availability}

Code, training logs including all failed runs, and the trained checkpoint are
released under Apache 2.0 at
\url{https://github.com/Laalinibh/protogrok-jax} and
\url{https://huggingface.co/Anvyon/protogrok-jax-306m}.

\bibliographystyle{plainnat}
\begin{thebibliography}{9}

\bibitem[Arp et al.(2022)]{arp2022dodos}
D.~Arp, E.~Quiring, F.~Pendlebury, A.~Warnecke, F.~Pierazzi, C.~Wressnegger,
L.~Cavallaro, and K.~Rieck.
\newblock Dos and don'ts of machine learning in computer security.
\newblock In \emph{USENIX Security Symposium}, 2022.

\bibitem[Chen et al.(2016)]{chen2016training}
T.~Chen, B.~Xu, C.~Zhang, and C.~Guestrin.
\newblock Training deep nets with sublinear memory cost.
\newblock \emph{arXiv:1604.06174}, 2016.

\bibitem[Friedman(2001)]{friedman2001greedy}
J.~H. Friedman.
\newblock Greedy function approximation: A gradient boosting machine.
\newblock \emph{Annals of Statistics}, 29(5):1189--1232, 2001.

\bibitem[Locatello et al.(2020)]{locatello2020object}
F.~Locatello, D.~Weissenborn, T.~Unterthiner, A.~Mahendran, G.~Heigold,
J.~Uszkoreit, A.~Dosovitskiy, and T.~Kipf.
\newblock Object-centric learning with slot attention.
\newblock In \emph{Advances in Neural Information Processing Systems}, 2020.

\bibitem[Loshchilov and Hutter(2019)]{loshchilov2019adamw}
I.~Loshchilov and F.~Hutter.
\newblock Decoupled weight decay regularization.
\newblock In \emph{International Conference on Learning Representations}, 2019.

\bibitem[Moustafa and Slay(2015)]{moustafa2015unsw}
N.~Moustafa and J.~Slay.
\newblock UNSW-NB15: A comprehensive data set for network intrusion detection
systems.
\newblock In \emph{Military Communications and Information Systems Conference
(MilCIS)}, 2015.

\bibitem[Pahalavan(2023)]{pahalavan2023nids}
R.~D. Pahalavan.
\newblock nids-datasets: Packet-level and flow-level data from UNSW-NB15 and
CIC-IDS2017.
\newblock \url{https://github.com/rdpahalavan/nids-datasets}, 2023.

\bibitem[Sommer and Paxson(2010)]{sommer2010outside}
R.~Sommer and V.~Paxson.
\newblock Outside the closed world: On using machine learning for network
intrusion detection.
\newblock In \emph{IEEE Symposium on Security and Privacy}, 2010.

\bibitem[Vaswani et al.(2017)]{vaswani2017attention}
A.~Vaswani, N.~Shazeer, N.~Parmar, J.~Uszkoreit, L.~Jones, A.~N. Gomez,
{\L}.~Kaiser, and I.~Polosukhin.
\newblock Attention is all you need.
\newblock In \emph{Advances in Neural Information Processing Systems}, 2017.

\end{thebibliography}

\end{document}
